\section{Affine events and private prime coordinates}\label{sec:affine-one-window}
A run event is a linear system over $\F$. Its probability is determined by
rank and compatibility; its dependence on a second event is determined by
relations between the two systems. This elementary observation lets us keep
arithmetic and probability separate.
For $n\ge1$, let
\[
 v(n)=(v_p(n)\bmod2)_p\in\F^{(\mathcal P)}.
\]
Write the prime signs additively: $X=(X_p)_{p\in\mathcal P}$
denotes the sequence of independent Bernoulli variables, uniform on
$\F$, defined by $f(p)=(-1)^{X_p}$. Then
\[
 f(n)=(-1)^{\langle X,v(n)\rangle}.
\]

For a start $x$, introduce the $L$ rows
\begin{align*}
 a_{x,0}&=v(x-1)+v(x), & b_{x,0}&=1,\\
 a_{x,j}&=v(x)+v(x+j), & b_{x,j}&=0,
 \qquad 1\le j<L.
\end{align*}
Let $A_x$ be the matrix with rows $a_{x,j}$ and $b_x=(1,0,\dots,0)$. Then
\[
 J_{x,L}=\1_{\{A_xX=b_x\}}.
\]
For a tuple $\mathbf x=(x_1,\dots,x_k)$, stack the matrices and the right-hand sides. Set
\[
 \calR(\mathbf x)=\ker(A_{\mathbf x}^{\mathsf T}),
 \qquad
 \rho(\mathbf x)=\dim\calR(\mathbf x),
\]
and
\[
 q_{\mathbf x}(c)=\langle c,b_{\mathbf x}\rangle.
\]

\begin{lemma}[Affine Fourier identity]\label{lem:fourier}
For every tuple $\mathbf x=(x_1,\dots,x_k)$,
\[
 2^{kL}\PP(J_{x_1,L}=\cdots=J_{x_k,L}=1)
 =\sum_{c\in\calR(\mathbf x)}(-1)^{q_{\mathbf x}(c)}.
\]
Consequently,
\[
 \PP(J_{x_1,L}=\cdots=J_{x_k,L}=1)
 =2^{-kL}\eta(\mathbf x)2^{\rho(\mathbf x)},
\]
where $\eta(\mathbf x)=1$ if $q_{\mathbf x}$ vanishes on $\calR(\mathbf x)$, and $0$ otherwise. Finally,
\begin{equation}\label{eq:affine-abs}
\left|\eta(\mathbf x)2^{\rho(\mathbf x)}-1\right|
 \le2^{\rho(\mathbf x)}-1.
\end{equation}
\end{lemma}

\begin{proof}
For a matrix with $m=kL$ rows,
\[
 \1_{\{AX=b\}}
 =2^{-m}\sum_{c\in\F^m}(-1)^{\langle c,AX-b\rangle}.
\]
Taking expectations, the terms with $A^{\mathsf T}c\ne0$ vanish. A character sum over a vector space equals the cardinality of that space if the character is trivial, and $0$ otherwise. If $\eta=0$, then $q$ is nontrivial, hence $\rho\ge1$, and $1\le2^\rho-1$.
\end{proof}


\subsection{The boundary map of a tree}
Put $B=L+1$, $I_L=\{-1,0,\ldots,L-1\}$ and
$V_x=\{x+i:i\in I_L\}$. The rows of $A_x$ are the edge vectors of the
star with edges $\{x-1,x\}$ and $\{x,x+j\}$, $1\le j<L$.
For a finite tree, the map taking an edge set to the vertices of odd degree
is a bijection onto the even vertex subsets. Thus a relation among the
rows of $A_x$ selects an even subset of $V_x$ whose product is a square.

\begin{lemma}[Square-product interpretation]\label{lem:even-subsets}
For disjoint windows $V_{x_1},\ldots,V_{x_k}$, the relation space is in
bijection with tuples of even subsets $S_i\subset V_{x_i}$ such that
$\prod_i\prod_{n\in S_i}n$ is a square. The affine character is
$\sum_i\1_{\{x_i-1\in S_i\}}$ modulo two.
More generally, the edge relations of any tree are identified with its
square-product even vertex subsets.
\end{lemma}
\begin{proof}
The sum of the selected edge vectors is the sum of the valuation vectors
at the odd-degree vertices. It vanishes exactly when their product is a
square. The boundary map is injective because a tree has no nonempty even
degree edge subgraph, and its image has dimension $|V|-1$, the dimension of
the even-subset space. The affine character records the root leaf edges.
\end{proof}

\subsection{One-window arithmetic}
For $H\ge2$ define the odd large-prime kernel
\[
 K_H(n)=\prod_{\substack{p>H\\v_p(n)\ \text{odd}}}p.
\]
Call $n$ \emph{$H$-defective} when $K_H(n)=1$. If an interval has diameter
less than $H$, every nondefective vertex has a prime coordinate appearing
there and nowhere else in the interval. Such a coordinate is a private
pivot. Notice that the definition uses odd valuations, not merely the
presence of a large prime divisor.

\begin{lemma}[Split-product Runge bound]\label{lem:runge}
There is an absolute $C_0$ such that the following holds. Let $d=2k\ge2$,
let $\gamma_1,\ldots,\gamma_d$ be distinct integers with
$|\gamma_i|\le R$, where $R\ge1$, and set $G(T)=\prod_i(T+\gamma_i)$.
If $U\in\Z$, $U\ge2R$, and $G(U)$ is an integer square, then
\begin{equation}\label{eq:runge-bound}
 U\le(C_0dR)^{2d}.
\end{equation}
\end{lemma}
\begin{proof}
Expand $\sqrt{G(T)}$ at infinity and let $P(T)$ be its polynomial part.
Its coefficients lie in $2^{-d}\Z$ and, for $U>4R$,
$|\sqrt{G(U)}-P(U)|\le(CdR)^d/U$.
If the difference is nonzero, integrality of $U$ and $\sqrt{G(U)}$ makes
its modulus at least $2^{-d}$. Otherwise the nonzero integer polynomial
$2^{2d}(G-P^2)$ vanishes at $U$; its coefficient height is at most
$(CdR)^{2d}$, and the integral-root bound applies. The range $2R\le U\le4R$
is immediate. All constants and the growing-degree coefficient estimates
are given in \cref{supp:lem:runge-calculation}. This is a direct
split-product implementation of Runge's method, not an unqualified
invocation of a fixed-degree theorem \cite{Runge1887,Walsh1992}.
\end{proof}

\begin{proposition}[Few defects, with their exponential weight]
\label{prop:defects}
Fix $0<c_1<c_2$. Uniformly for integers $c_1\log X\le H\le c_2\log X$, an
interval of $H+1$ integers in $[X/2,3X]$ contains
$O_{c_1,c_2}(\log X/\log\log X)$ $H$-defective integers. If $m_u$ counts
these defects in $[u,u+H]$, then
\begin{equation}\label{eq:defect-weight}
 \sum_{X\le u<2X}(2^{m_u}-1)\le X^{1/2+o_{c_1,c_2}(1)}.
\end{equation}
\end{proposition}
\begin{proof}
Write each defective integer as $sa^2$, where $s$ is squarefree and
$H$-smooth. Associate to the $m$ defects in a window the columns
$((v_p(s))_{p\le H},1)\in\F^{\pi(H)+1}$. A nonzero kernel word selects
an even number $d$ of distinct integers with square product. The Runge
bound gives $X/2\le(CdH)^{2d}$; consequently its Hamming weight is at
least $c\log X/\log H$ for a fixed $c>0$.
Let $t$ be an integer of order $\log X/\log H$ with
$2t+1\le c\log X/\log H$. If the kernel is zero, $m\le\pi(H)+1$.
If $m<2t$, the desired bound is again immediate. Otherwise disjoint
Hamming balls about its words give
\[
 \sum_{j\le t}\binom mj\le2^{\pi(H)+1},\qquad
 (m/t)^t\le2^{\pi(H)+1}.
\]
Since $(\pi(H)+1)/t=O_{c_1,c_2}(1)$, it follows that
$m=O_{c_1,c_2}(\log X/\log\log X)$.

On the other hand the number of defective integers up to $3X$ is at most
\[
 \sqrt{3X}\sum_{\substack{s\ \text{squarefree}\\P^+(s)\le H}}s^{-1/2}
 =\sqrt{3X}\prod_{p\le H}(1+p^{-1/2})=X^{1/2+o(1)}.
\]
Each belongs to at most $H+1=X^{o(1)}$ windows. Only
$X^{1/2+o(1)}$ windows therefore have a defect, and the pointwise bound
makes each weight $2^{m_u}$ equal to $X^{o(1)}$. This proves the sum.
\end{proof}

For $0<\delta<1$, write $U_{M,\delta}=[\lceil M^\delta\rceil,M)\cap\Z$.
If $B\asymp\log M$, a dyadic slice $X\le x<2X$ in this interval has
$B\asymp_\delta\log X$. Applying \cref{prop:defects} on these slices gives
\begin{equation}\label{eq:macro-defects}
 \sum_{x\in U_{M,\delta}}(2^{m_B(x)}-1)\le M^{1/2+o(1)},\qquad
 m_B(x)=\#\{n\in V_x:K_B(n)=1\}.
\end{equation}
The same assertion holds for windows whose width and defect threshold are
any two comparable quantities in a fixed logarithmic band: use an enclosing
window of width equal to a fixed multiple of the threshold and subdivide
it into a bounded number of windows before the coding argument. In the
applications below the threshold can always be chosen at least the width.

\begin{corollary}[Uniform first moment and an affine bound]
\label{cor:first-moment}
Put $p=2^{-L}$ and $m_B^\circ(x)=\#\{n\in V_x\setminus\{x-1\}:K_B(n)=1\}$.
For any right-hand side $b\in\F^L$,
\begin{equation}\label{eq:one-affine-error}
 \left|\PP(A_xX=b)-p\right|
 \le p(2^{m_B^\circ(x)}-1),\qquad
 \PP(A_xX=b)\le2^{-L+\max(m_B(x)-1,0)}.
\end{equation}
For every deterministic $A\subset I_N$ and $B$ in a fixed logarithmic band,
\[
 \EE\sum_{x\in A}J_{x,L}=|A|p+O(pN^{1/2+o(1)}).
\]
The analogous error over any subset of $U_{M,\delta}$ is $O(pM^{1/2+o(1)})$.
\end{corollary}
\begin{proof}
A private pivot excludes its vertex from any relation. An even subset of
$m$ defective vertices has at most $\max(m-1,0)$ degrees of freedom.
Thus $\rho(x)\le\max(m_B(x)-1,0)\le m_B^\circ(x)$.
Apply \cref{lem:fourier}, which is unchanged by the affine right-hand side,
and sum \cref{prop:defects} or \eqref{eq:macro-defects}. Restricting a
nonnegative error sum to a mask can only decrease it.
\end{proof}


For $Y\ge2$ write
$\mathcal F_Y=\sigma(f(p):p\text{ prime},\ p\le Y)$ for the information
carried by the small prime signs.

\begin{corollary}[Prescribed values in one window]\label{cor:absolute-marginal}
Let $M_x$ have rows $v(n)$, $n\in V_x$, and let $w\in\{\pm1\}^B$.
Write $I_{x,w}=\1_{\{(f(x-1+j))_{0\le j<B}=w\}}$ and $p_B=2^{-B}$.
Then
\begin{equation}\label{eq:absolute-marginal}
 |\PP(I_{x,w}=1)-p_B|\le p_B(2^{m_B(x)}-1).
\end{equation}
If, for every vertex $n\in V_x$, there is a prime $p>Y$, where $Y>B$,
such that $v_p(n)$ is odd, then
$\PP(I_{x,w}=1\mid\mathcal F_Y)=p_B$ for every realization of
$\mathcal F_Y$.
Fix $0<c_1<c_2$. Uniformly for $c_1\log N\le B\le c_2\log N$,
every deterministic start mask $A\subset I_N$ and dictionary
$\mathcal W\subset\{\pm1\}^B$ satisfy
\[
 \EE\sum_{x\in A}\sum_{w\in\mathcal W}I_{x,w}
 =|A|\,|\mathcal W|\,2^{-B}
 +O_{c_1,c_2}\!\left(|\mathcal W|\,2^{-B}N^{1/2+o_{c_1,c_2}(1)}\right).
\]
The expectation is unconditional. The threshold and error are uniform in
$B,A,\mathcal W$, which may vary with $N$; empty masks and dictionaries
are allowed. The mask restricts starts, not vertices: every occurrence
uses all $B$ values $x-1,\ldots,x+B-2$, without truncation at the mask or
at the boundary of $I_N$.
\end{corollary}
\begin{proof}
For a fixed window of positive integers, take a prime cutoff at least
as large as its largest vertex, and at least $Y$ when conditioning.
The word event in the infinite product
is the inverse image of the corresponding finite-cylinder event, so their
probabilities agree exactly. Only the prime signs are independent; the
integer values retain all their multiplicative relations.
In a relation among the rows of $M_x$, every nondefective vertex is
excluded by its private coordinate. Hence the nullity is at most $m_B(x)$,
without the even-subset restriction used for relative signs. Finite Fourier
inversion gives \eqref{eq:absolute-marginal}. On a good support, the
projection of $M_x$ onto primes above $Y$ has an identity minor of order
$B$, so every prescribed value vector has probability $2^{-B}$ on each
positive-mass assignment atom of $\mathcal F_Y$. Sum the nonnegative
bound using \cref{prop:defects}, then use linearity of expectation.
The same threshold works for every mask and every dictionary; no
independence between occurrences is needed.
\end{proof}

\begin{lemma}[Private pivots on a graph]\label{lem:private-tree}
Let $G$ be a finite graph. In each connected component suppose all vertices
except possibly one have distinct private prime coordinates in the whole
graph. The kernel of its edge vectors, projected onto these coordinates,
is contained in its cycle space. In particular its rank is at least
$|E(G)|-\beta_1(G)$, where $\beta_1(G)=|E(G)|-|V(G)|+c(G)$.
For a tree this rank is $|E(G)|$.
\end{lemma}
\begin{proof}
A vanishing edge combination has even degree at every pivoted vertex.
The sum of degrees in a connected component is even, so its possible
unpivoted vertex has even degree as well. The selected edge set is therefore
in the cycle space, of dimension $\beta_1(G)$.
\end{proof}

\begin{lemma}[Local run pairs]\label{lem:local-pairs}
If $0<|x-y|<L$, then $J_{x,L}J_{y,L}=0$. Over ordered pairs with
$|x-y|=L$ in a dyadic block or in $U_{M,\delta}$, the homogeneous weights
$2^{\rho(x,y)}$ have sum at most $N^{1+o(1)}$ or $M^{1+o(1)}$, respectively.
\end{lemma}
\begin{proof}
An overlapping start requires two consecutive signs within the first
constant run to be opposite. At distance $L$, the two start trees form a
tree on an interval of diameter $2L$. With threshold $2L+1$, private
pivots and \cref{prop:defects} bound its relation dimension by
$O(\log N/\log\log N)$ on a slice. There are $O(N)$ such pairs;
summing slices gives the macroscopic assertion.
\end{proof}
